A satellite orbits the Earth on a circular orbit. The initial momentum of the
satellite is given by the vector $\mathbf{p}$. At a certain time, an explosive
charge is set off which gives the satellite an additional impulse
$\Delta\mathbf{p}$, equal in magnitude to $|\mathbf{p}|$. Let $\alpha$ be the
angle between the vectors $\mathbf{p}$ and $\Delta\mathbf{p}$, and $\beta$
between the radius vector of the satellite and the vector $\Delta\mathbf{p}$.
By thinking about the direction of the additional impulse $\Delta\mathbf{p}$,
consider if it is possible to change the orbit to each of the cases given
below. If it is possible mark YES on the answer sheet and give values of
$\alpha$ and $\beta$ for which it is possible. If the orbit is not possible,
mark NO.

\begin{description}
\item[(a)] a hyperbola with perigee at the location of the explosion.
\item[(b)] a parabola with perigee at the location of the explosion.
\item[(c)] an ellipse with perigee at the location of the explosion.
\item[(d)] a circle.
\item[(e)] an ellipse with apogee at the location of the explosion.
\end{description}

Note that for $\alpha = 180^\circ$ and $\beta = 90^\circ$ the new orbit will be
a line along which the satellite will free fall vertically towards the centre
of the Earth.
